
\subsubsection{2.1 Curriculum Learning}

Bengio et al. (2009) introduced curriculum learning: training on examples ordered from easy to hard. Hacohen \& Weinshall (2019) demonstrated benefits for image classification. Platanios et al. (2019) applied curriculum principles to neural machine translation. These methods operate at the example level using heuristics like loss-based difficulty.

The approach in this paper differs by operating at the domain level rather than example difficulty, and by using corpus statistics rather than training-time metrics.

\subsubsection{2.2 Multi-Domain Pretraining}

GPT-3, PaLM, and Llama train on diverse corpora spanning web text, books, code, and scientific papers. Standard practice uses static domain mixing: fixed sampling probabilities (e.g., Common Crawl 67\%, Books 16\%, Wikipedia 4.5\%) maintained throughout training.

Xie et al. (2023) introduced DoReMi, which optimizes static mixing ratios using group distributionally robust optimization. DoReMi demonstrated that tuned static mixtures outperform uniform mixing by 2-6\% across benchmarks. DoReMi's weights remain static throughout training.

The approach in this paper introduces temporal dynamics. While DoReMi determines optimal static proportions, diversity-ranked scheduling orders temporal presentation given fixed total exposure per domain.

\subsubsection{2.3 Multi-Phase Training}

Two-phase training is common: general pretraining on broad corpora followed by continued training on domain-specific data. This represents a coarse temporal curriculum with a single sharp transition.

Chronopoulou et al. (2019) found that task ordering affects final performance in multi-task learning. Ruder \& Plank (2017) studied temporal aspects of data selection for neural sequence labeling. These works focus on task-level transfer in supervised settings.

The approach in this paper extends multi-phase training to smooth, parametric curriculum schedules. Instead of binary phase transitions, it implements continuous Gaussian-weighted domain scheduling.

\subsubsection{2.4 Gradient Geometry}

Stringer et al. (2019) introduced the participation ratio to measure gradient covariance rank in neuroscience models. Fort et al. (2019) demonstrated that later layers exhibit higher gradient diversity during training.

Kornblith et al. (2019) proposed Centered Kernel Alignment (CKA) for measuring representational similarity across neural networks. Neyshabur et al. (2020) argued that early training shapes the geometry of the solution basin.

This paper proposes using these tools to test a mechanistic explanation: early high-diversity training may establish broader gradient covariance (measurable via Participation Ratio), which persists through path-dependent optimization (measurable via CKA similarity between early and final checkpoints). This mechanism is not verified in the current proof-of-concept.

\subsubsection{2.5 Continual Learning}

Kirkpatrick et al. (2017) introduced Elastic Weight Consolidation, which penalizes changes to parameters important for previous tasks. Rebuffi et al. (2017) proposed rehearsal methods that replay previous data during new task learning.

McCloskey \& Cohen (1989) identified catastrophic forgetting: learning new information impairs performance on previously learned tasks. Ramasesh et al. (2021) studied forgetting in large language models during continued pretraining.

The hypothesis (not tested in this work) is that pretraining-time intervention through diversity-ranked scheduling may establish broad gradient geometry that reduces catastrophic forgetting during later continual learning.
